\documentclass[UTF8,11pt]{ctexart}
\usepackage{amsmath,amssymb,mathtools,geometry}
\geometry{a4paper,margin=20mm}\allowdisplaybreaks
\begin{document}
\section*{分时电价下电动公交排班与满充联合模型}

\paragraph{来源} PDF第3--5页，§2.2，式(1)--(13)。按正文解释明确条件弧与时间索引。
\subsection*{集合、参数与变量}
$V$ 为有固定开始时刻 $a_i$、结束时刻 $b_i$ 及任务耗电 $u_i$ 的公交班次，
$R$ 为单一充电站的访问副本，$o,e$ 为同一车场的出发和返回副本。
$A$ 包含允许的 $o\to V$、$V\to V$、$V\to R$、$R\to V$、$V\to e$ 弧；
不含自环、充电站间弧。副本数须足够每个班次后一次充电使用。
$t_{ij},d_{ij}$ 为弧旅行时间和耗电；电池容量为 $Q$，车辆日固定成本 $c_f$，
充电速率 $g>0$（电量/时间），车数上限 $N_{\rm EB}$。
$h_i$ 为班次 $i$ 结束后至充电站或车场所需能耗的原文安全界。
$E(t)$ 是式(1)积分中的分时\emph{充电成本率}；若输入为电量单价 $p(t)$，
应先提供 $E(t)=g p(t)$ 以使积分单位为货币。原文表2将其称电价，但式(1)直接按时间积分，
费用按各时段电价及上述单位换算计算。

$x_{ij}\in\{0,1\}$ 表示连接弧；$q_i\in[0,Q]$ 表示到达班次/充电站时的电量；
$\tau^a_{ij},\tau^b_{ij}\ge0$ 为班次 $i$ 后到充电副本 $j$ 的开始/结束时刻。
未选择的充电弧两时刻固定为零，其成本也为零。
\subsection*{完整核心模型}
\[
\min\ c_f\sum_{j\in V}x_{oj}
 +\sum_{i\in V,j\in R:(i,j)\in A}
  x_{ij}\int_{\tau^a_{ij}}^{\tau^b_{ij}}E(t)\,dt.\tag{1}
\]
\begin{align}
\sum_{j:(i,j)\in A}x_{ij}&=1 &&i\in V,\tag{2}\\
\sum_{j:(i,j)\in A}x_{ij}&\le1 &&i\in R,\tag{3}\\
\sum_{i:(i,j)\in A}x_{ij}&=\sum_{k:(j,k)\in A}x_{jk}
&&j\in V\cup R,\tag{4}\\
\sum_{j\in V}x_{oj}&=\sum_{i\in V}x_{ie}\le N_{\rm EB}.\tag{12}
\end{align}
如下时间/能量关系只作用于实际选择的弧；这将原文大M激活逻辑写成无歧义的条件形式：
\begin{align}
x_{ij}=1&\Longrightarrow b_i+t_{ij}\le a_j &&i,j\in V,\tag{5}\\
x_{ij}=1&\Longrightarrow \tau^a_{ij}=b_i+t_{ij} &&i\in V,\ j\in R,\tag{6}\\
x_{ij}=1&\Longrightarrow \tau^b_{ij}=\tau^a_{ij}+(Q-q_j)/g
&&i\in V,\ j\in R,\tag{7}\\
x_{ij}=x_{jk}=1&\Longrightarrow \tau^b_{ij}+t_{jk}\le a_k
&&i,k\in V,\ j\in R,\tag{8}\\
x_{ij}=1&\Longrightarrow q_j\le q_i-u_i-d_{ij} &&i,j\in V,\tag{9a}\\
q_j&\ge u_j+h_j &&j\in V,\tag{9b}\\
x_{ij}=1&\Longrightarrow 0\le q_j\le q_i-u_i-d_{ij}
&&i\in V,\ j\in R,\tag{10}\\
x_{oj}=1&\Longrightarrow 0\le q_j\le Q-d_{oj} &&j\in V,\tag{11a}\\
x_{ij}=1&\Longrightarrow 0\le q_j\le Q-d_{ij} &&i\in R,\ j\in V,\tag{11b}\\
x_{ie}=1&\Longrightarrow q_i-u_i-d_{ie}\ge0 &&i\in V,\tag{11c}\\
x_{ij}&\in\{0,1\},\quad 0\le q_i\le Q,\quad
0\le\tau^a_{ij}\le\tau^b_{ij}.\tag{13}
\end{align}
式(11c)明确回场能量，若 $h_i$ 已按原文覆盖车场耗电则它冗余。
只允许从 $o$ 到 $e$ 的连通班次链；固定班次的严格时间顺序及正旅行/充电时间排除环，
退化零时长输入需加入标准连通性割，禁止无仓闭环。
\paragraph{原式与整理说明}
式(5)原下标印成 $j\in R$，但该式及其文字是两个连续班次，改为 $j\in V$；
式(8)在充电弧及其后继弧同时选中时生效。
原(9)--(10)未选弧的 $Q$ 激活项仍残留 $u_i$，这里按正文“连续班次/充电弧”
用条件关系澄清，避免未选弧错误约束电量。电量状态采用原式给出的上界关系。
原旧LaTeX中的逆向积分和乘以 $g$ 已依PDF恢复为开始至结束积分与除以 $g$。

\end{document}
